\documentclass[11pt,a4paper]{article}
\usepackage[margin=24mm,headheight=15pt]{geometry}
\usepackage{fontspec}
\setmainfont{Times New Roman}
\setsansfont{Avenir Next}
\setmonofont{Menlo}
\usepackage{microtype}
\usepackage{newtxmath}
\usepackage{mathtools}
\usepackage{booktabs,longtable,array,tabularx}
\usepackage{pdflscape}
\usepackage[table]{xcolor}
\usepackage{enumitem}
\usepackage{fancyhdr}
\usepackage{titlesec}
\usepackage{xurl}
\usepackage{hyperref}
\usepackage{bookmark}
\usepackage{fancyvrb}
\usepackage[most]{tcolorbox}

\definecolor{Navy}{HTML}{18324A}
\definecolor{Accent}{HTML}{C56A3A}
\definecolor{LinkBlue}{HTML}{2D6A8A}
\definecolor{SoftBlue}{HTML}{EAF1F5}
\definecolor{TableHead}{HTML}{DCE8EF}
\definecolor{BodyGray}{HTML}{26343D}

\hypersetup{
  colorlinks=true,
  linkcolor=LinkBlue,
  urlcolor=LinkBlue,
  citecolor=LinkBlue,
  pdftitle={Finite certificate for the four-star layered-poset construction},
  pdfauthor={Final Thesis Master Project}
}
\setlength{\parindent}{0pt}
\setlength{\parskip}{0.65em}
\setlength{\emergencystretch}{3em}
\setlist{leftmargin=1.6em,itemsep=0.3em,topsep=0.25em}
\setcounter{tocdepth}{2}
\renewcommand{\arraystretch}{1.15}
\allowdisplaybreaks

\titleformat{\section}
  {\Large\bfseries\sffamily\color{Navy}}
  {\thesection}{0.65em}{}
  [\vspace{-0.25em}\color{Accent}\titlerule]
\titleformat{\subsection}
  {\large\bfseries\sffamily\color{Navy}}
  {\thesubsection}{0.6em}{}
\titleformat{\subsubsection}
  {\normalsize\bfseries\sffamily\color{LinkBlue}}
  {\thesubsubsection}{0.55em}{}

\newtcolorbox{ResearchQuote}{
  enhanced,
  breakable,
  colback=SoftBlue,
  colframe=LinkBlue,
  boxrule=0pt,
  leftrule=2.2pt,
  arc=0pt,
  left=8pt,right=8pt,top=5pt,bottom=5pt
}
\newenvironment{ResearchContents}
  {\begin{tcolorbox}[breakable,colback=SoftBlue,colframe=SoftBlue,arc=2pt]}
  {\end{tcolorbox}}

\pagestyle{fancy}
\fancyhf{}
\fancyhead[L]{\small\sffamily\color{Navy} Asymmetric stochastic TSP}
\fancyhead[R]{\small\sffamily\color{Navy}\detokenize{four-star-poset-finite-certificate.md}}
\fancyfoot[C]{\small\sffamily\color{BodyGray}\thepage}
\renewcommand{\headrulewidth}{0.3pt}
\renewcommand{\footrulewidth}{0pt}

\begin{document}
\color{BodyGray}
\begin{titlepage}
  \thispagestyle{empty}
  \vspace*{22mm}
  {\sffamily\small\bfseries\color{Accent}\MakeUppercase{Research note}}\par
  \vspace{8mm}
  {\sffamily\fontsize{20}{24}\selectfont\bfseries\color{Navy} Finite certificate for the four-star layered-poset construction}\par
  \vspace{6mm}
  {\color{Accent}\rule{42mm}{1.4pt}}\par
  \vspace{11mm}
  {\Large\color{BodyGray} Asymmetric stochastic TSP\\
  clairvoyance-gap investigation}\par
  \vfill
  \begin{tcolorbox}[
    colback=SoftBlue,colframe=SoftBlue,arc=2pt,
    left=10pt,right=10pt,top=8pt,bottom=8pt]
    \sffamily
    \textbf{Source}\quad \detokenize{four-star-poset-finite-certificate.md}\\[3pt]
    \textbf{Prepared}\quad 25 July 2026\\[3pt]
    \textbf{Format}\quad Typeset mathematical PDF
  \end{tcolorbox}
  \vspace{12mm}
\end{titlepage}

\begin{ResearchContents}
\tableofcontents
\end{ResearchContents}
\clearpage

This note records one completely finite member of the four-star layered-poset family and checks, with slack, that its asymmetric clairvoyance gap is strictly larger than \(4/3\).  It also states precisely what the asymptotic argument proves: a liminf of at least \(3/2\), rather than an equality unless a separate adaptive upper bound is supplied.


The construction and the causal product lemma are proved in \nolinkurl{four-star-layered-poset-construction.md}.  Here we use only the two bounds proved there.  With \(L\) stages, selector probability \(q\), and


\[
                    h=q^2(1-q)^2,
\]

they are


\begin{equation}
\begin{aligned}
 \inf_\pi \mathbb E K_\pi
   &\ge 3-(1-h)^L,\\
 \mathbb E W
   &\le 2+L(4q^3-2q^4).
\end{aligned}
\tag{2}
\end{equation}

Here \(K_\pi\) is the number of increasing runs in the service order of a causal policy and \(W\) is the width of the realized active subposet.


\phantomsection
\section*{1. Concrete parameters}
\addcontentsline{toc}{section}{1. Concrete parameters}

Take


\begin{equation}
             q=\frac1{200},\qquad L=56000,\qquad
             \varepsilon=10^{-10}.
\tag{3}
\end{equation}

There are \(2(L+1)=112002\) permanent gate clients and \(4L=224000\) stochastic selector clients, hence \(336002\) clients in total.


For these values,


\[
 h=\frac{39601}{1600000000}
\]

and


\begin{equation}
 Lh
 =\frac{56000\cdot39601}{1600000000}
 =1.386035.
\tag{4}
\end{equation}

Therefore


\begin{equation}
 (1-h)^L\le e^{-Lh}<0.26.
\tag{5}
\end{equation}

The final numerical inequality has ample slack.  For example, the first five nonzero terms of the exponential series already give


\[
 e^{1.386035}
 >
 1+1.386035+\frac{1.386035^2}{2}
 +\frac{1.386035^3}{6}
 +\frac{1.386035^4}{24}
 >\frac{50}{13},
\]

which is equivalent to the strict part of (5).


It follows from (1) that


\begin{equation}
                 \inf_\pi\mathbb E K_\pi>2.74.
\tag{6}
\end{equation}

For the posterior bound,


\[
 4q^3-2q^4
 =\frac{399}{800000000},
\]

so (2) gives


\begin{equation}
 \mathbb EW
 \le 2+\frac{56000\cdot399}{800000000}
 =2.02793.
\tag{7}
\end{equation}

\phantomsection
\section*{2. Conversion to a finite directed metric}
\addcontentsline{toc}{section}{2. Conversion to a finite directed metric}

For distinct clients \(x,y\), put


\begin{equation}
 d(x,y)=
 \begin{cases}
   \varepsilon,&x<y\text{ in the layered poset},\\
   1,&x\not<y,
 \end{cases}
\qquad
 d(r,x)=d(x,r)=\frac12.
\tag{8}
\end{equation}

This is a positive directed metric.  Two consecutive \(\varepsilon\)-arcs imply \(x<y<z\), hence \(x<z\), and every other two-client path has length at least the direct distance.  A path through the depot has length one and merely ties an incomparable client arc.


If \(N\) clients are active and their service order has \(K\) maximal increasing runs, its closed-tour cost is exactly


\begin{equation}
              \varepsilon N+(1-\varepsilon)K.
\tag{9}
\end{equation}

The expected number of active clients is


\begin{equation}
 \mu:=\mathbb EN
 =2(L+1)+4Lq
 =113122.
\tag{10}
\end{equation}

Consequently (6)--(10) imply


\begin{equation}
\begin{aligned}
 \operatorname{OPT}_{\rm adapt}
  &>\varepsilon\mu+(1-\varepsilon)2.74,\\
 \operatorname{OPT}_{\rm post}
  &\le\varepsilon\mu+(1-\varepsilon)2.02793.
\end{aligned}
\tag{11}
\end{equation}

Multiplying the desired strict inequality by three and using (11),


\begin{equation}
\begin{aligned}
 3\operatorname{OPT}_{\rm adapt}
 -4\operatorname{OPT}_{\rm post}
 &>
 (1-\varepsilon)
   \bigl(3\cdot2.74-4\cdot2.02793\bigr)
 -\varepsilon\mu\\
 &=(1-\varepsilon)0.10828-0.0000113122\\
 &>0.10826>0.
\end{aligned}
\tag{12}
\end{equation}

Thus this one finite product-distribution instance satisfies


\[
 \boxed{\displaystyle
 \frac{\operatorname{OPT}_{\rm adapt}}
      {\operatorname{OPT}_{\rm post}}>\frac43.}
\]

No parallel repetition, correlated activation, zero distance, or small-instance computation is used.


\phantomsection
\section*{3. Exact asymptotic conclusion}
\addcontentsline{toc}{section}{3. Exact asymptotic conclusion}

The parameter choice


\[
 q_L=L^{-2/5},\qquad \varepsilon_L=L^{-2}
\]

in the main construction gives


\[
 \operatorname{OPT}_{\rm adapt}\ge3-o(1),
 \qquad
 \operatorname{OPT}_{\rm post}=2+o(1).
\]

Hence the proved statement is


\begin{equation}
 \liminf_{L\to\infty}
 \frac{\operatorname{OPT}_{\rm adapt}}
      {\operatorname{OPT}_{\rm post}}
 \ge\frac32.
\tag{13}
\end{equation}

Equation (13) is stronger than the required strict \(4/3\) lower bound. Calling it a limit equal to \(3/2\) would additionally require an adaptive upper bound of \(3+o(1)\), which is not needed for the construction theorem.

\end{document}
